\chapter*{Introduction}
\addcontentsline{toc}{chapter}{Introduction}

The basic question with which this book is concerned is: 'What are the politics of Marxism?' It is not a simple question, and a preliminary discussion of the reasons why it is not simple may be the best way to begin answering it.

The first difficulty lies in the term 'Marxism' itself. Marx died in 1883 and his own 'Marxism' (a term, incidentally, which he never used) was subsequently greatly enlarged, first of all by Engels in the twelve very active and prolific years by which he survived Marx; and then by a number of other major and not so major figures in the following years up to the present. The most prominent of these figures by far was Lenin, whose own successors invented the term 'Marxism-Leninism', not only to denote Lenin's contribution to the enlargement of Marxism, but also to proclaim as settled what was in fact a highly contentious question, namely how firm was the link between the Marxism of Marx on the one hand and that of Lenin on the other. The question is particularly contentious, as it happens, in regard to some central aspects of the politics of Marxism, for instance the role assigned to the party.

In one way or another, the same problem arises in relation to all the major figures who, by word or deed, have left their stamp on Marxism after Marx. The issue is by no means 'academic'.

Given the political importance of Marxism in the twentieth century and its use as a political weapon, it is on the contrary a question which has itself strong and even explosive political implications and consequences.

A second difficulty in the discussion of Marxism and politics has to do with the character of the writings on politics of all the major figures of Marxism, beginning with Marx himself. These writings are for the most part the product of particular historical episodes and specific circumstances; and what there is of theoretical exploration of politics in what may be called classical Marxism (by which I mainly mean the writings of Marx, Engels, and Lenin, and, at a different level, those of some other figures such as Rosa Luxemburg, Gramsci, and Trotsky) is mostly unsystematic and fragmentary, and often part of other work. Some of the most basic texts of the politics of Marxism 2 Marxism and Politics answer to this description, for instance Marx's Eighteenth Brumaire of Louis Bonaparte and The Civil War in France, or Lenin's What is to be Done? and The State and Revolution. In fact, there are very few classical Marxist political texts which do not answer to this description.* This is not meant to devalue the significance and interest of these works but only to note that none of the greatest figures of classical Marxism, with the partial exception of Gramsci, ever attempted or for that matter felt the need to attempt the writing of a 'political treatise'. Given their total engagement in political struggles, and the vital importance they all attached to theory, this is very remarkable and cannot be taken as accidental. It has in fact to do with the meaning and status of 'politics ' in Marxism and requires separate discussion later in this chapter. For the moment, it may be noted that this absence of systematic political theorization on the part of Marx, Engels, and their most prominent successors means in effect that a Marxist politics has to be constructed or reconstructed from the mass of variegated and fragmented material which forms the corpus of Marxism. The danger which this presents of arbitrary selection and emphasis is obvious, and is further enhanced by a strong tendency to extremely one-sided interpretations which Marxism engenders, not least among Marxists. This danger is more easily acknowledged than avoided; but one must do the best one can.

At the same time—and this is another difficulty to bear in mind—it is as well to recognize that what I have just called the 'corpus of Marxism' has very definite limitations in terms of the construction or reconstruction of a Marxist politics. One of these limitations is that the available classical writings are simply silent or extremely perfunctory over major issues of politics and political theory: there is a limit to what can properly be squeezed out of a paragraph, a phrase, an allusion or a metaphor. The point is particularly obvious in regard to a whole range of political experience in the last fifty years or so, which is of crucial importance in any attempt at theorizing the politics of Marxism, but about which classical Marxism is, in any direct way, naturally silent. After all, Marx did die in 1883 and Engels in 1895, * In the short Preface which he wrote in 1869 for the second edition of The Eighteenth Brumaire, Marx noted that the work had been written in the form of weekly articles and 'under the immediate pressure of events' (SE, p. 142).

1 his could be said of most of the relevant political texts of classical Marxism.

Introduction 3 Rosa Luxemburg in 1919, and Lenin in 1924. Gramsci was in gaol from 1926 until shortly before his death in 1937, and wonderfully suggestive though his Prison Notebooks are, it has to be recognized that they are no more than fragments and notes, written under dismal conditions and in the grip of disease. Only Trotsky, until his assassination in 1940, provides a sustained commentary, out of classical Marxism, on a world that encompassed Stalinism and Fascism and much else as well; and this too was written under exceptional pressure and in very special circumstances.

It might well be asked why more by way of a Marxist political theorization of some of the most important experiences of our times, and of a Marxist political theory in general, should not have been constructed in, say, the last fifty years, on the foundations provided by classical Marxism.

At least a part of the answer must be sought in the experience of Stalinism and of Stalinism's domination of Marxism over a period of some thirty years and more from the late twenties onwards. The point is not so much that the Stalinist version of Marxism was a dreadfully impoverished affair, though it was certainly that; but that its version of Marxism, and style of approach to it, were accepted, largely because of extraneous political reasons, by the vast majority of people who then called themselves Marxists, not only in Russia but everywhere else as well.

A particular quality of Stalinism, in this realm as in all others, was its imperative and binding definition of 'the line' to be followed. This served to decide, in catechismal form, what were the 'fundamentals ' of Marxism, or rather of 'Marxism-Leninism '; and it also served to specify who was and—even more important—who was not to be regarded as having made a contribution to Marxist thought. Those who were not so regarded constituted a great host—in fact most of the people who have in this century made a serious contribution to the development of Marxism. So they were banned in the Soviet Union and greatly neglected or altogether ignored by most Marxists outside.

Marx and Engels were not banned. But they were read very selectively; and probably less than their accredited interpreters, Lenin and Stalin. The accent was on authoritative interpretations and non-arguable propositions; and on the capacity of 'Marxism-Leninism'and dialectical materialism (another term 4 Marxism and Politics which Marx never used or even knew) to resolve all theoretical problems, or at the very least to provide a sure guide to their solution. These were not perspectives in which Marxist thought could be expected to flourish, least of all in the highly sensitive area of political theory and political analysis; and it did not.

Things have changed in the last twenty-odd years, at least in those countries where Marxism is not the official ideology, and even up to a point in those where it is. The spell was at long last broken in the fifties; and these two decades have been a period of intense probing and questioning in Marxist thought which is without parallel in its history. But there is a vast amount of ground to make up, most of all perhaps in the area of Marxist politics.

The probing and questioning which have occurred have been fed from many sources. The first of them is the very fact that Stalin's regime came under attack in the Soviet Union itself: however superficial and inadequate that attack was and has remained since the XXth Party Congress of 1956, it did release a vast array of key questions which can never again be ignored or suppressed. There has also been the experience of China and the many challenges which it has posed to the Russian and East European 'models ' of socialism. These challenges remain even if one rejects the attempts which have been made to create a Chinese shrine at which to worship in place of the discredited Russian one.

There have also been the liberation struggles all over the 'Third World', and their reverberations in the countries of advanced capitalism. In these countries, class struggles which had in the post-war years been declared to belong to the past have brought back on the agenda the largest questions concerning reform and revolution; and these struggles have been given an added dimension by women's movements, ethnic and student movements, and new forms of expression and action. The inchoate and many-sided 'cultural revolution' which has gripped advanced capitalist countries in these two decades, has also and naturally affected Marxism and forced many reassessments upon it.

Of course, these stirrings within Marxism have their own limitations, particularly in the countries where it is the official ideology. It is in this latter respect a grim and sobering thought that in the infinite number of words produced in Communist Introduction 5 countries in the last two decades, not one article could safely appear which, for instance, resolutely attacked Lenin on any issue of importance. As for capitalist countries, the recovery of Marxism has involved the proliferation of coteries and fashionable trends, with different sects and schools proclaiming their own version of Marxism as the only authentic one, or as the only one appropriate to the times we live in, or whatever. But this is the inevitable price which a live and vigorous movement pays for its liveliness and vigour. However it may be in China, Chairman Mao's old injunction to let a hundred flowers bloom has been taken at face value in capitalist countries; and it is of no great account that a fair number of weeds should in the process have grown as well.

The important change, in this realm, is that there is now no 'recognized' Marxist orthodoxy to which general obedience is paid, except where such obedience is capable of being imposed; and this is very different from what happened in the past, when such an orthodoxy was internationally accepted, and not even perceived as an orthodoxy. The change is a great gain. It was after all Marx who once said that his favourite motto was 'Doubt all things'; and it is all to the good that there should be less and less people in favour of amending this to read: 'Doubt all things except what Comrade X, Chairman Y, or President Z pronounces upon.' There are many worse slogans than 'Everyone his own Marx'. For in the end, there is no 'authoritative' interpretation—only personal judgement and evaluation.

In relation to the politics of Marxism, this requires that the first textual priority and attention should be given to Marx himself, and to Engels. This is the essential starting-point, and the only possible 'foundation' of Marxism-as-politics. It is only after this has been done that one can usefully take up Lenin, Luxemburg, Gramsci, and others to try and construct a Marxist politics.

It is worth saying at the outset, however, that the most careful textual scrutiny will not yield a smooth, harmonious, consistent, and unproblematic Marxist political theory. On the contrary, it is only by a very superficial reading, or by fiat, that such a Marxist political theory, of a distorted kind, can be obtained.

Not only are the texts susceptible to different and contradictory interpretations: they also do actually incorporate tensions, contradictions, and unresolved problems which form an intrinsic 6 Marxism and Politics part of Marxist political thought. To ignore or to try and obscure this does not simply distort the real nature of that thought but deprives it of much of its interest.

I noted earlier how remarkable it was that no major Marxist figure had tried to set out systematically the substance and specificity of Marxist political theory. The people concerned were after all among the most gifted and penetrating minds of the last hundred years, men and women utterly immersed in political life, struggle, and ideas; and the same people have always attached the highest importance to theory as an indispensable part of class struggle and working-class politics.

'Without theory, no revolutionary movement' is one of the precepts of Lenin which Marxists have most readily accepted; and they have understood it to mean that without a clear articulation of its theoretical premisses and projections, a working-class movement advances blindly. This makes the absence in classical Marxism of a systematic theorization of Marxist politics all the more remarkable. As I have already suggested, the reasons for this must be sought in the concept of politics which informs Marxism, and they are in fact deeply embedded in the structure of Marxist thought concerning social life and the place of politics in it.

On the most general plane, Marxism begins with an insistence that the separation between the political, economic, social, and cultural parts of the social whole is artificial and arbitrary, so that, for instance, the notion of 'economics ' as free from 'politics ' is an ideological abstraction and distortion. There is no such thing as 'economics'—only 'political economy', in which the 'political' element is an ever-present component.

On this view, politics is the pervasive and ubiquitous articulation of social conflict and particularly of class conflict, and enters into all social relations, however these may be designated. But this very pervasiveness of politics appears to rob it of its specific character and seems to make it less susceptible to particular treatment, save in the purely formal description of processes and institutions which Marxists have precisely wanted to avoid. In reality, it is perfectly possible to treat politics as a specific phenomenon, namely as the ways and means whereby social conflict and notably class conflict is manifested.

At one end, this may mean accommodation and agreement between social groups which are not greatly divided (or for that Introduction 7 matter which are); at the other end, it may mean civil war which, to adapt Clausewitz, is politics carried out by other means.

There is, however, a more particular and direct reason for the Marxist neglect of political theory, which has to do with the concept of 'base' and 'superstructure', or rather with the implications which have tended to be drawn from it.

One of the most influential formulations in the whole body of Marxist thdught occurs in a famous passage of the 'Preface' to A Contribution to the Critique of Political Economy of 1859 and goes as follows: In the social production of their life, men enter into definite relations that are indispensable and independent of their will, relations of production which correspond to a definite stage of development of their material productive forces. The sum total of these relations of production constitutes the economic structure of society, the real foundation on which rises a legal and political superstructure and to which correspond definite forms of social consciousness. The mode of production of material life conditions the social, political and intellectual life process in general. It is not the consciousness of men that determines their being, but, on the contrary, their social being that determines their consciousness.1 Another text of Marx, which makes the same point in a somewhat different form, is also worth quoting, given the importance of this whole conceptualization for the status of the political element in Marxism. In the third volume of Capital, Marx writes the specific economic form, in which unpaid surplus-labour is pumped out of direct producers, determines the relationship of rulers and ruled, as it grows directly out of production itself and, in turn, reacts upon it as a determining element ... it is always the direct relationship of the owners ol the conditions of production to the direct producers—a relation always naturally corresponding to a definite stage in the development of the methods of labour and thereby its social productivity—which reveals the innermost secret, the hidden basis of the entire social structure, and with it the political form of the relation of sovereignty and dependence, in short, the corresponding specific form of the state.2 Clearly, these texts can easily be interpreted as turning politics into a very 'determined' and 'conditioned' activity indeed—so determined' and 'conditioned', in fact, as to give politics a mostly derivative, subsidiary, and 'epiphenomenaT character. At its extreme, this turns Marxism into an 'economic determinism' which deprives politics of any substantial degree of autonomy.

8 Marxism and Politics For their part, Marx and Engels explicitly rejected any rigid and mechanistic notion of 'determination'; and Engels specifically repudiated the idea that Marx and he had ever intended to suggest that 'the economic element is the only determining one which he described as a 'meaningless, abstract, senseless phrase'.3 As for Marx, the passage from Capital just quoted goes on to say that 'the same economic basis', because of 'innumerable different empirical circumstances, natural environment, racial relations, external historical influences, etc.', will show 'infinite variations and gradations in appearance, which can be ascertained only by analysis of the empirically given circumstances';4 and these 'variations' must obviously include the political part of the 'superstructure'.

'Base' and 'superstructure' must be taken, in Gramsci's phrase, as the elements of an 'historical bloc'; and the different elements which make up that 'bloc' will vary in their relative weight and importance according to time, place, circumstance, and human intervention.

But one must not protest too much. There remains in Marxism an insistence on the 'primacy' of the 'economic base' which must not be understated. This 'primacy' is usually taken by Marxists, following Engels,5 as meaning that the 'economic base ' is decisive, or determining, 'in the last instance'. But it is much more apposite and meaningful to treat the 'economic base ' as a starting-point, as a matter of the first instance. In a different but relevant context, Marx noted that in all forms of society there is one specific kind of production which predominates over the rest, whose relations thus assign rank and influence to the others. It is a general illumination which bathes all the other colours and modifies their particularity. It is a particular ether which determines the specific gravity of every being which has materialized within it.6 This formulation may well be applied to the relation which politics, in the Marxist scheme, has to the 'economic base '; and it is then possible to proceed from there, and to attribute to political forms and forces whatever degree of autonomy is judged in any particular case to be appropriate. In this usage, the notion of 'primacy' constitutes an important and illuminating guideline, not an analytical straitjacket. The ways in which that 'primacy' determines and conditions political and other forms remain to be discovered, and must be treated in each case as specific, Introduction 9 circumstantial, and contingent; and this also leaves open for assessment the ways in which political forms and processes in turn affect, determine, condition, and shape the economic realm, as of course they do and as they are acknowledged to do by Marxists, beginning with Marx.

Marx's own cast of mind was strongly anti-determinist, and led him to reject all trans-historical and absolute determinations, beginning with the Hegelian 'determination ' of the historical process, and including the rejection of any attempt, as he wrote with heavy sarcasm in 1877, to use 'as one's master key a general historico-philosophical theory, the supreme virtue of which consists in being super-historical'.7 Marx did believe that certain things must come to pass, notably the supersession of capitalism: but a belief in the inevitability of certain events is not the same as a belief in their particular 'determination'.

The question then is not whether Marxism is an 'economic determinism'. I take it that it is not. The point is rather that the entirely legitimate emphasis which Marxists have placed on the importance of the economic 'infra-structure' and the mode of production has resulted, in relation to social analysis and notwithstanding ritual denegations concerning 'economic determinism', in a marked 'economisin' in Marxist thought.

The term 'economisin' has now come to be used in very loose ways and has been made to cover a multitude of sins, real or imaginary.* But in the present context, it means both the attribution of an exaggerated—almost an exclusive—importance to the economic sphere in the shaping of social and political relations, leading precisely to 'economic determinism'; and it also involves a related underestimation of the importance of the 'superstructural' sphere. In the letter of Engels to Bloch quoted earlier, he also wrote that 'Marx and I are ourselves partly to blame for the fact that the younger people sometimes lay more stress on the economic side than is due to it. We had to *In recent usage, it has meant (a) the belief that the public ownership of the means of production can be equated with, or is at least bound to be followed by, the socialist transformation of the 'relations of production (b) the belief that a massive development of the productive forces is an essential precondition for the achievement of socialist 'relations of production'; and (c) that with the abolition of capitalist owners, the state altogether changes its character and comes to reflect or incarnate the 'dictatorship of the proletariat'. How far (or whether (these propositions constitute deformations of Marxism is a matter for discussion.

10 Marxism and Politics emphasise the main principle, vis-a-vis our adversaries, who denied it, and we had not always the time, the place or the opportunity to allow the other elements involved in the interaction to come into their rights.,H It may be noted that Engels did not here renounce the 'main principle he only regretted the fact that it had sometimes been allowed to obscure or crowd out the 'other elements5.

However, not all the elements of the 'superstructure' suffered equal neglect by Marxists. In intellectual matters, classical Marxism was deeply concerned with economic analysis but also with history and philosophy, and other areas of thought as well, for instance science. It was political theory which, comparatively speaking, suffered the greatest neglect; and this has remained so, even in the more recent decades of Marxist intellectual growth.

The reason for this may be traced back to a fundamental distinction which Marx drew' at the very beginning of his political life between 'political emancipation' and 'human emancipation'. Political emancipation, by which he meant the achievement of civic rights, the extension of the suffrage, representative institutions, the curbing of monarchical rule, and the curtailment of arbitrary state power in general, was by no means to be scorned. On the contrary, it should be welcomed, Marx wrote in his essay On the Jewish Question of 1843, as 'certainly a big step forward' and as 'the last form of human emancipation within the prevailing scheme of things V The stress is Marx s own and is significant: it points to a major Marxist theme, namely that human emancipation can never be achieved in the political realm alone but requires the revolutionary transformation of the economic and social order. 'Property, etc., in short the whole content of law and the state', Marx also wrote in his Critique of Hegel's 'Philosophy of Right' in the same year, 'is broadly the same in North America as in Prussia. Hence the republic in America is just as much a mere form of the state as the monarchy here. The content of the state lies beyond these constitutions.'10 There was great strength in the Marxist insistence that sense could not be made out of political reality without probing beneath political institutions and forms; and that insistence was and remains the basis of Marxist political analysis and of Marxist political sociology.

Introduction 11 But while this in no way requires the conclusion that political forms within the prevailing scheme of things', are of no great consequence, it produces a tendency to draw just such a conclusion; and the tendency has been very strong within Marxism to devalue or ignore the importance of 'mere ' political forms and to make very little of the problems associated with them.

This tendency was further reinforced by an extraordinarily complacent view of the ease with which political problems (other than mastering bourgeois resistance) would be resolved in post-revolutionary societies. Politics was taken to be an expression of man's alienation. 'Human emancipation' meant, among other things, the end of politics. As Istvan Meszaros has summarized it, politics must be conceived as an activity whose ultimate end is its own annulment by means of fulfilling its determinate function as a necessary stage in the complex process of positive transcendence.'11 In The Poverty of Philosophy (1846), Marx wrote that 'the working class, in the course of its development, will substitute for the old civil society* an association which will exclude classes and their antagonism, and there will be no more political power properly so-called, since political power is precisely the expression of antagonism in civil society. ,12Lucio Colletti has rightly noted14 that this was one of the most consistent themes of Marx throughout his life, from the Critique of Hegel's Philosophy of Right' to The Civil War in France of 1871 and the Critique of the Gotha Programme of 1875. It is instructive in this connection to see how contemptuously (and inadequately) Marx answered the very pertinent questions which Bakunin, in his Statism and Anarchism of 1874, had raised about some of the problems which he thought must arise in the attempt to bring about the rule of the proletariat, in the literal sense in which Marx meant it.14 But it was Engels who gave to the notion of the end of politics its most popular Marxist expression in his Anti- Diihring of 1878, the relevant part of which was reproduced in Socialism: Utopian and Scientific, published in 1892 and probably, after the Communist Manifesto, the most widely-read of the works of Marx and Engels. Writing about the state, Engels said that * 'Civil society ' here stands for bourgeois society. See K. Marx and F. Engels, The German Ideology (London, 1965), p. 48.

12 Marxism and Politics when at last it becomes the real representative of the whole of society, it renders itself unnecessary. As soon as there is no longer any social class to be held in subjection; as soon as class rule, and the individual struggle for existence based upon our present anarchy of production, with the collisions and excesses arising from these, are removed, nothing more remains to be repressed, and a special repressive force, a state, is no longer necessary . . . State interference in social relations becomes, in one domain after another, superfluous, and then withers away of itself; the government of persons is replaced by the administration of things, and by the conduct of processes of production. The state is not 'abolished'. It withers away.15 This optimistic view was reaffirmed in the most extreme form in Lenin's The State and Revolution, which was written on the eve of the October Revolution of 1917, and where all the problems of the exercise of socialist power—for instance the danger of the bureaucratization of the revolution and the reproduction of a strongly hierarchical order, not to mention the question of civic freedoms—were either swept aside or simply ignored. This complacency was replaced almost as soon as the Bolsheviks had seized power by a sombre awareness on the part of Lenin and some of the other Bolshevik leaders of how genuine and difficult these problems were, and how real was the threat that they posed to the new Soviet order. But it is significant that neither Lenin nor those around him should have thought it necessary at least to consider these problems seriously in the years preceding the revolution. They did take up many theoretical and organizational problems in the years before 1917, from economic analysis and philosophy to the organization of the party. But save for the debate among Marxists around the issue of the relation of the party to the working class and the danger of 'substitutism', which followed the publication of Lenin's What is to Be Done? in 1902, there was comparatively little attention devoted to the theoretical and practical problems posed by the concept of socialist democracy and the 'dictatorship of the proletariat'. It is symptomatic that it was only after 1917 that one of the most alert minds among the Bolsheviks, Bukharin, should have become aware of the import of the challenge posed by theories of elite and of bureaucratization which had been posed to Marxists—and left unanswered for a good many years past.16 Much of the reason for this must be sought in the absence in Marxism of a serious tradition of political inquiry; and in the assumption commonly made by Marxists before 1917 that the Introduction 13 socialist revolution would itself—given the kind of overwhelming popular movement it would be—resolve the main political problems presented to it. The State and Revolution was the ultimate expression of that belief.

There did occur very extensive and passionate discussion of all such problems after the Russian Revolution, in Russia and in the socialist movements of many other countries; and this would undoubtedly have produced in time a strong Marxist tradition of political thought. But the debates were overtaken by Stalinist triumphalism' from the mid-twenties onwards; and this required agreement that the main political problems of socialist power had been solved in the Soviet Union, except for the problem presented by the enemies of socialism. To suggest otherwise was to cast doubt on the socialist and proletarian character of Soviet rule, and thus to turn automatically into one of the very enemies of the Soviet Union and socialism. The impact of such modes of thought on the history of Marxism in subsequent years cannot be overestimated.

There is one altogether different explanation for the relative poverty of Marxist political theory which has been advanced over the years, namely and very simply that, whatever the contribution of Marxism may be in other areas of thought, it has little to contribute in this one. This is what anti-Marxist political theorists and others have commonly held to be the case: Marxist politics and political pronouncements might be an object of study, but not much more.

Perhaps more remarkably, a fairly negative verdict on Marxist political theory has also been returned in recent years by Lucio Colletti, who is one of the most interesting modern Marxist writers. In his Introduction to Marx's Early Writings which I have already mentioned, Colletti states that Marx, in his Critique of Hegel's 'Philosophy of Right', 'already possesses a very mature theory of politics and the state'.'The Critique, after all', he goes on, 'contains a clear statement of the dependence of the state upon society, a critical analysis of parliamentarism accompanied by a counter-theory of popular delegation, and a perspective showing the need for the ultimate suppression of the state itself. He then concludes that 'politically speaking, mature Marxism would have relatively little to add to this.'17 Colletti also says of Lenin's The State and Revolution that it 'advances little beyond the ideas set out in the Critique'; and 14 Marxism and Politics even though he refers to Marx's 'profundity ' in connection with the Critique; and to the 'marvellous continuation' in On the Jewish Question of the 'theory of the state elaborated in the Critique, he takes the view that 'Marxism's most specific terrain of development was the socio-economic one. ,18For Colletti, the most important progenitor of Marx's political theory was Rousseau; and 'so far as “political” theory in the strict sense is concerned, Marx and Lenin have added nothing to Rousseau, except for the analysis (which is of course rather important) of the “economic bases” for the withering away of the state'.19 Claims such as these are too broad and qualified to offer a firm basis of argument, and there does not seem much point in entering into abstract disputation as to whether the originality of Marxism lies in one realm rather than in another. The same applies to the judgements of anti-Marxist writers on the worth of Marxist political theory. The best way to test all such propositions is to show what a Marxist political theory specifically involves; and to indicate how far it may serve to illuminate any particular aspect of historical or contemporary reality.

For this purpose, the developments in Marxist political thinking in recent years have obviously been of great value, not least because the constricting 'triumphalism' of an earlier period has been strongly challenged; and the challenge has produced a much greater awareness among Marxists that Marxism, in this as much as in other realms, is full of questions to be asked and—no less important—of answers to be questioned. Many hitherto neglected or underestimated problems have attracted greater attention; and many old problems have been perceived in a better light. As a result, the beginnings have been made of a political theorization in the Marxist mode.

But these are only beginnings, and in some areas barely even that. One such area is that covered by Communist experience since 1917 and the nature and workings of Communist states and political systems. There has not really been very much, beyond Trotsky's The Revolution Betrayed of forty years ago, by way of Marxist attempts to theorize the experience of Stalinism; and the Marxist debate on the nature of the Soviet state (and of other Communist regimes) has long been paralysed by the invocation of formulas and slogans—'degenerate workers' state' versus 'state capitalism' and so forth. The whole area has largely been left for anti-Marxists to explore; and the subject Introduction 15 badly requires serious and sustained Marxist political analysis and reinterpretation. This has now begun but needs to be pushed much further.

The position is in many ways rather better in regard to the politics of the many different countries which are arbitrarily subsumed under the label Third World'. But here too, it would appear to the non-specialist that, as far as political analysis is concerned, no more than some paths have been cleared, and that the main work of theorizing the known practice remains to be undertaken; and it is only in the undertaking of it that it will be possible to discover which theoretical categories of Marxism are relevant to the experience in question, which need to be modified, and which should be discarded.

Perhaps not surprisingly, it is the countries of advanced capitalism which have received most attention from Marxists in the last two decades, in the area of politics as well as in other areas of inquiry. Even so, it is not the wealth of Marxist political theory and analysis in relation to these societies which is striking, but the fact that it has not proceeded much further. Nothing like enough serious work has yet been done which could be said to constitute a Marxist tradition of political studies.

The present work attempts to make a contribution to the development of such a tradition. It suggests some of the main questions which must be considered and probed in the construction of a Marxist political analysis; and it does so on the basis of a reading of primary Marxist texts, first and foremost on the basis of a reading of Marx himself. Whether the reading is accurate or not must be left for others to judge. So too must the question of the validity of the political argument which, in this reconstruction of Marxist politics, I am inevitably led to put forward and which indeed I want to put forward.

NOTES 1 K. Marx, 'Preface' to A Contribution to the Critique of Political Economy, in SW 1968, p. 182.

2 K. Marx, Capital (Moscow, 1962), III, p. 772.

3 F. Engels to J. Bloch, 21-2 September 1890, in SW 1968, p. 692.

4 K. Marx, Capital, III, p. 772.

5 See e.g. his letter to C. Schmidt, 27 October 1890, in SW 1968, pp. 694-9.

6 K. Marx, Grundrisse (London, 1973), pp. 106-7.

7 SC. p. 379. The letter was to the Editorial Board of the journal Otechestven- niye Zapiski and is dated November 1877.

8 SW 1968, p. 692.

16 Marxism and Politics 9 On the Jewish Question, in EW, p. 221.

10 K. Marx, Critique of Hegel's Doctrine of the State, ibid., p. 89. The usual title of the work is that given in the text above.

11 I. Meszaros, Marx's Theory of Alienation (London, 1970), p. 160.

12 K. Marx, The Poverty of Philosophy (London, 1946), p.147.

13 See L. Colletti's Introduction to EW.

14 For an extract from Marx's comments on the book, seeFI (London, 1974). For the full text, see Marx-Engels, Werke (Berlin, 1964), XVIII.

15 F. Engels, Anti-Diihring (Moscow, 1962), p. 385.

16 See S. F. Cohen, Bukharin and the Bolshevik Revolution (London, 1974), p.

21.

17 EW, p. 45.

18 Ibid., p. 46.

19 L. Colletti, 'Rousseau as Critic of “Civil Society'”, in From Rousseau to Lenin (London, 1972), p. 185.